% !TeX root = ../main.tex
\section{Tail Latency}\label{sec:tail}
\rubric{Tail latency and device state}{25}{p50/p95/p99/p99.9 reporting and interpretation of queueing,
thermal and garbage-collection effects (with \cref{sec:device-state}).}

\begin{hypothesis}
  Moving from a mid queue depth to the near-knee depth should raise the median modestly but stretch the
  upper percentiles much more: deeper queues mean more requests waiting behind the occasional slow
  flash operation, so the p99/p99.9 gap over p50 should widen.
\end{hypothesis}

\begin{knobs}
  Fixed       & \qty{4}{\kibi\byte} random read; one job; \qty{6}{\second} runs, 5 repetitions; NVMe
                temperature recorded per repetition \\
  Varied      & iodepth 8 (mid), 32 and 64 (bracketing the knee of \cref{sec:queue-depth}), 128 (past
                it) \\
  Script      & \code{SSD_Performance_Profiling/scripts/run.sh tail} \\
\end{knobs}

\begin{figure}[tbp]
  \centering
  \IfDataExists{tail}{%
    \begin{tikzpicture}
      \begin{axis}[acs, ybar, bar width=7pt, width=0.78\linewidth, ymode=log, log origin=infty,
                   symbolic x coords={8,32,64,128}, xtick=data, enlarge x limits=0.2,
                   xlabel={iodepth}, ylabel={Latency (\si{\micro\second})}, legend pos=north west, legend columns=4]
        \addplot+[iqr bars] table[x=qd, y=lat_p50_us_median,  iqr of=lat_p50_us]  {data/tail.csv};  \addlegendentry{p50}
        \addplot+[iqr bars] table[x=qd, y=lat_p95_us_median,  iqr of=lat_p95_us]  {data/tail.csv};  \addlegendentry{p95}
        \addplot+[iqr bars] table[x=qd, y=lat_p99_us_median,  iqr of=lat_p99_us]  {data/tail.csv};  \addlegendentry{p99}
        \addplot+[iqr bars] table[x=qd, y=lat_p999_us_median, iqr of=lat_p999_us] {data/tail.csv};  \addlegendentry{p99.9}
      \end{axis}
    \end{tikzpicture}%
  }
  \caption{Latency percentiles at a mid queue depth (8), around the knee (32, 64) and past it (128);
  medians with IQR bars. These runs followed \ensuremath{\sim}\qty{11}{\gibi\byte} of random writes
  while the drive was still doing background flash work, and the upper tail detaches from the median as
  QD grows: p99.9/p50 rises from about 4$\times$ to 24$\times$.}\label{fig:tail}
\end{figure}

\subsection{Analysis}
The tail runs did \emph{not} see the same drive as the queue-depth read sweep, and that difference is the
main result. In \code{run.sh all} they started immediately after the queue-depth \emph{write} sweep,
which wrote \ensuremath{\sim}\qty{11}{\gibi\byte} of random \qty{4}{\kibi\byte} data, and there was no
idle gap. At the same queue depth, the tail runs delivered half the IOPS of the earlier, clean read
sweep: \csvq{tail}{qd=32}{iops_median} vs.\ \csvq{queue_depth}{qd=32,mode=read}{iops_median} at QD=32.
Their median latency was
\csvq[0]{tail}{qd=32}{lat_p50_us_median} vs.\ \csvq[0]{queue_depth}{qd=32,mode=read}{lat_p50_us_median}\,\si{\micro\second}.
The temperature was \qtyrange{50.9}{52.9}{\celsius}, so this is not thermal. The likely cause is
foreground reads competing with the controller's background folding and garbage collection of the data
just written.

Under that contention the tail detaches from the body as QD grows, as hypothesised but more strongly:
\begin{itemize}[nosep]
  \item QD=8: p50~\csvq{tail}{qd=8}{lat_p50_us_median}, p99~\csvq{tail}{qd=8}{lat_p99_us_median},
        p99.9~\csvq{tail}{qd=8}{lat_p999_us_median}\,\si{\micro\second}; p99.9/p50 $=4.1\times$.
  \item QD=32: p50~\csvq{tail}{qd=32}{lat_p50_us_median}, p99~\csvq{tail}{qd=32}{lat_p99_us_median},
        p99.9~\csvq{tail}{qd=32}{lat_p999_us_median}\,\si{\micro\second}; $6.7\times$.
  \item QD=64: p99.9 \csvq{tail}{qd=64}{lat_p999_us_median}\,\si{\micro\second}; $15\times$.
  \item QD=128: p99.9 \csvq{tail}{qd=128}{lat_p999_us_median}\,\si{\micro\second}; $24\times$.
\end{itemize}
Meanwhile IOPS barely grow past QD=32 (\csvq{tail}{qd=32}{iops_median}$\to$\csvq{tail}{qd=128}{iops_median}).
In the clean sweep, by contrast, p99.9/p50 stays between 3.8 and $4.9\times$ from QD=16 to 256: the whole
distribution shifts up together, which is pure queueing. The tail ratio is therefore a sensitive
detector of background flash activity. Queueing alone keeps it near $4\times$, and a read that arrives
during a program/erase or a garbage-collection burst waits milliseconds. Deeper queues hold more requests
behind each such stall.

These are the numbers that govern provisioning for a latency-SLA service, not the mean. At QD=32 on a
drive that has recently absorbed writes, one read in a thousand takes
\csvq{tail}{qd=32}{lat_p999_us_median}\,\si{\micro\second}. The knee from \cref{sec:queue-depth}
is therefore a latency decision as much as a throughput one: past it, almost no IOPS are gained and these
tails inflate fastest. The NVMe composite temperature stayed at \qtyrange{33.9}{61.9}{\celsius} across
all runs, far below the \qty{86.85}{\celsius} warning threshold (\cref{sec:device-state}).
